\textbf{Phần 1. Trắc nghiệm lựa chọn}
\setcounterpageref{ex}{0}

\begin{ex} %[4D1N04]
Biểu thức nào sau đây là biểu thức chứa một chữ?
\choice
{$15 + 24$}
{\True $35 + a$}
{$a + b$}
{$a + b + c$}
\loigiai{
$35 + a$ là biểu thức chứa một chữ $a$.\\
\textbf{Đáp án của câu hỏi là B}
}
\end{ex}

\begin{ex} %[4D1N04]
Biểu thức $a + b$ là biểu thức chứa
\choice
{\True Hai chữ}
{Một chữ}
{Ba chữ}
{Không chứa chữ}
\loigiai{
$a + b$ chứa 2 chữ $a$ và $b$.\\
\textbf{Đáp án của câu hỏi là A}
}
\end{ex}

\begin{ex} %[4D1H04]
Giá trị của biểu thức $45 + a$ với $a = 15$ là
\choice
{$50$}
{$55$}
{$65$}
{\True $60$}
\loigiai{
$45 + 15 = 60$.\\
\textbf{Đáp án của câu hỏi là D}
}
\end{ex}

\begin{ex} %[4D1H04]
Giá trị của biểu thức $a - b$ với $a = 80, b = 25$ là
\choice
{$45$}
{$65$}
{\True $55$}
{$50$}
\loigiai{
$80 - 25 = 55$.\\
\textbf{Đáp án của câu hỏi là C}
}
\end{ex}

\begin{ex} %[4D1H04]
Giá trị của biểu thức $m \times n$ với $m = 12, n = 5$ là
\choice
{$50$}
{$70$}
{$55$}
{\True $60$}
\loigiai{
$12 \times 5 = 60$.\\
\textbf{Đáp án của câu hỏi là D}
}
\end{ex}

\begin{ex} %[4D1H04]
Biểu thức tính chu vi $P$ của hình vuông có độ dài cạnh là $a$ là
\choice
{$P = a + 4$}
{$P = a \times a$}
{\True $P = a \times 4$}
{$P = a : 4$}
\loigiai{
Chu vi hình vuông là $P = a \times 4$.\\
\textbf{Đáp án của câu hỏi là C}
}
\end{ex}

\begin{ex} %[4D1V04]
Biểu thức tính chu vi hình chữ nhật có chiều dài $a$ và chiều rộng $b$ (cùng đơn vị đo) là
\choice
{\True $P = (a + b) \times 2$}
{$P = a + b \times 2$}
{$P = a \times b$}
{$P = (a + b) : 2$}
\loigiai{
$P = (a + b) \times 2$.\\
\textbf{Đáp án của câu hỏi là A}
}
\end{ex}

\begin{ex} %[4D1V04]
Giá trị của biểu thức $a + b + c$ với $a = 24, b = 36, c = 40$ là
\choice
{$90$}
{$110$}
{\True $100$}
{$95$}
\loigiai{
$24 + 36 + 40 = 60 + 40 = 100$.\\
\textbf{Đáp án của câu hỏi là C}
}
\end{ex}

\begin{ex} %[4D1V04]
Một hình chữ nhật có $a = 15\text{ cm}, b = 8\text{ cm}$. Diện tích $S = a \times b$ của hình chữ nhật đó là
\choice
{$110\text{ cm}^2$}
{$130\text{ cm}^2$}
{$46\text{ cm}^2$}
{\True $120\text{ cm}^2$}
\loigiai{
$S = 15 \times 8 = 120\text{ cm}^2$.\\
\textbf{Đáp án của câu hỏi là D}
}
\end{ex}

\begin{ex} %[4D1V04]
Với $m = 48, n = 6$, giá trị của biểu thức $m : n + 15$ là
\choice
{$21$}
{\True $23$}
{$25$}
{$27$}
\loigiai{
$48 : 6 + 15 = 8 + 15 = 23$.\\
\textbf{Đáp án của câu hỏi là B}
}
\end{ex}

\begin{ex} %[4D1V04]
Cho biểu thức $A = 120 - (x + y)$. Với $x = 35, y = 45$, giá trị của $A$ là
\choice
{\True $40$}
{$30$}
{$50$}
{$60$}
\loigiai{
$A = 120 - (35 + 45) = 120 - 80 = 40$.\\
\textbf{Đáp án của câu hỏi là A}
}
\end{ex}

\begin{ex} %[4D1C04]
Cho biểu thức $P = a \times 5 + b \times 5$. Với $a + b = 24$, giá trị của biểu thức $P$ là
\choice
{$100$}
{\True $120$}
{$140$}
{$150$}
\loigiai{
$P = 5 \times (a + b) = 5 \times 24 = 120$.\\
\textbf{Đáp án của câu hỏi là B}
}
\end{ex}

\textbf{Phần 2. Trắc nghiệm đúng sai}
\setcounterpageref{ex}{0}

\begin{ex} %[4D1H04]
Xét tính đúng sai khi tính giá trị biểu thức:
\choiceTF
{\True Với $a = 5$, giá trị của $20 + a \times 3 = 35$.}
{Với $a = 10$, giá trị của $(45 - a) : 5 = 6$.}
{\True Với $a = 10$, giá trị của $(45 - a) : 5 = 7$.}
{\True Với $x = 8, y = 4$, giá trị của $x \times y : 2 = 16$.}
\loigiai{
\begin{itemize}
    \item $20 + 5 \times 3 = 20 + 15 = 35$. Đúng.
    \item $(45 - 10) : 5 = 35 : 5 = 7 \neq 6$. Ý 2 Sai, Ý 3 Đúng.
    \item $8 \times 4 : 2 = 32 : 2 = 16$. Đúng.
\end{itemize}
\begin{itemize}
    \item a) \textbf{Đúng}.
    \item b) \textbf{Sai}.
    \item c) \textbf{Đúng}.
    \item d) \textbf{Đúng}.
\end{itemize}
}
\end{ex}

\begin{ex} %[4D1V04]
Cho biểu thức $M = a + b - c$. Xét tính đúng sai:
\choiceTF
{\True $M$ là biểu thức chứa ba chữ $a, b, c$.}
{\True Khi $a = 50, b = 30, c = 20$ thì $M = 60$.}
{Khi $a = 100, b = 40, c = 50$ thì $M = 80$.}
{\True Khi $a = 100, b = 40, c = 50$ thì $M = 90$.}
\loigiai{
\begin{itemize}
    \item $M$ chứa 3 chữ. Đúng.
    \item $50 + 30 - 20 = 60$. Đúng.
    \item $100 + 40 - 50 = 90 \neq 80$. Ý 3 Sai, Ý 4 Đúng.
\end{itemize}
\begin{itemize}
    \item a) \textbf{Đúng}.
    \item b) \textbf{Đúng}.
    \item c) \textbf{Sai}.
    \item d) \textbf{Đúng}.
\end{itemize}
}
\end{ex}

\textbf{Phần 3. Trả lời ngắn}
\setcounterpageref{ex}{0}

\begin{ex} %[4D1H04]
Tính giá trị của biểu thức $150 - a$ với $a = 65$.
\shortans{85}
\loigiai{
\textbf{Đáp số:} $85$\\
$150 - 65 = 85$.
}
\end{ex}

\begin{ex} %[4D1H04]
Tính giá trị của biểu thức $a \times 8 + 24$ với $a = 7$.
\shortans{80}
\loigiai{
\textbf{Đáp số:} $80$\\
$7 \times 8 + 24 = 56 + 24 = 80$.
}
\end{ex}

\begin{ex} %[4D1V04]
Tính chu vi của hình vuông có cạnh $a = 18\text{ cm}$.
\shortans{72}
\loigiai{
\textbf{Đáp số:} $72$\\
$18 \times 4 = 72\text{ cm}$.
}
\end{ex}

\begin{ex} %[4D1V04]
Tính giá trị của biểu thức $a + b \times c$ với $a = 25, b = 15, c = 4$.
\shortans{85}
\loigiai{
\textbf{Đáp số:} $85$\\
$25 + 15 \times 4 = 25 + 60 = 85$.
}
\end{ex}

\begin{ex} %[4D1V04]
Tính chu vi hình chữ nhật với chiều dài $a = 24\text{ cm}$, chiều rộng $b = 16\text{ cm}$.
\shortans{80}
\loigiai{
\textbf{Đáp số:} $80$\\
$(24 + 16) \times 2 = 80\text{ cm}$.
}
\end{ex}

\begin{ex} %[4D1C04]
Cho biểu thức $S = a + b + c + d$. Biết $a + b = 50$ và $c + d = 75$. Tính giá trị của biểu thức $S$.
\shortans{125}
\loigiai{
\textbf{Đáp số:} $125$\\
$S = (a + b) + (c + d) = 50 + 75 = 125$.
}
\end{ex}
